%%% ВЫБОР УСТРОЙСТВА %%%

\IfStrEqCase{\devicecode}{%
{ЮТКБ.656122.611-06.01 РЭ}{%ДЗТ
% код для этого типа
\def\TypeTr{НН1(НН)} 
\def\fileExt{t}
\def\AbzazOne{
\item
Функция защиты от потери охлаждения контролирует уровни тока в обмотках ВН, НН1(НН) и НН2(СН). К каждой из обмоток подключено по два реле тока для контроля за уровнем нагрузки трансформатора. Каждое реле имеет индивидуально настраиваемый порог по току срабатывания, задаваемый пользователем.
}
}%
{ЮТКБ.656122.611-06.02 РЭ}{%ДЗАТ
% код для этого типа
\def\TypeTr{НН} 
\def\fileExt{at}
\def\AbzazOne{
\item
Функция защиты от потери охлаждения контролирует уровни тока в обмотках ВН, НН, а также в общей обмотке автотрансформатора. К каждой из обмоток подключено по два реле тока для контроля за уровнем нагрузки автотрансформатора. Каждое реле имеет индивидуально настраиваемый порог по току срабатывания, задаваемый пользователем.
}
}%
}
[
\def\TypeTr{??}
\def\fileExt{??}
\def\AbzazOne{??}
]


\color{unimain}{\subsection{Защита от потери охлаждения}}
\color{black}

\begin{enumerate}

\item
Защита от потери охлаждения направлена на предотвращение перегрева и повреждения первичного оборудования (автотрансформатора, шунтирующего ректора) в результате частичного или полного отказа системы охлаждения.

\AbzazOne

\item
Защита содержит три ступени, каждая из которых может быть выполнена с контролем температуры масла и уровня тока. Ввод контроля перегрева масла, а также выбор соответствующего уровня тока осуществляется программными накладками \textbf{«ЗПО-N конт.ДТм»} и \textbf{«ЗПО-N конт.РТ»} соответственно (где N – номер ступени). Каждая из ступеней имеет индивидуальную регулируемую задержку на срабатывание.

\item
Алгоритм работы ЗПО приведен на рисунке \ref{pic:ZPO}. Параметры для настройки ЗПО приведены в таблице \ref{app_set:zpo}.

\medskip
\begin{figure}[H]
\centering
\includegraphics[width=1\textwidth,height=1\textheight,keepaspectratio]{\fbpath/06.02 ДЗАТ/041.1151 ЗПО/img/ZPO_\fileExt.pdf}
\captionof{figure}{Функциональная схема логики ЗПО}\label{pic:ZPO}
\end{figure}

\end{enumerate}